\chapter{Programming Interface for a C++ Client: libCBview}
\label{cha:libcbview}

The libCBview provides a C++ encapsulation of libCB.
It provides only an object-oriented API for ConceptBase
and does not provide any additional methods in contrast to
libCB.

Compilation and linking has to be done in the same way as for
libCB. Note that you have to link both libraries libCB and libCBview
if you want to use the C++ classes.

The documentation in the following sections has been generated with
DOC++ (\url{http://docpp.sourceforge.net}).
% man kann mit CB_Make docpptex ein TeX-File generieren, dass die Dokumentation enth�lt.
% Dieses File wurde aber nachtraeglich editiert, daher sollte man Aenderungen besser
% manuell hier nachtragen.

\begin{cxxclass}
{class}
        {CBclient}
        {}
        {A client class for ConceptBase }
        {2}
\begin{cxxfunction}
{}
        {CBclient}
        {()}
        {Constructs an "empty" client which is not connected}
        {2.1}
\begin{cxxdoc}
Constructs an "empty" client which is not connected
\end{cxxdoc}
\end{cxxfunction}
\begin{cxxfunction}
{}
        {CBclient}
        {(char*\ host,\ int\ port,\ char*\ tool=(char*)NULL,\ char*\ user=(char*)NULL)}
        { Constructs a new CBclient object and connect to the specified host}
        {2.2}
\begin{cxxdoc}

Constructs a new CBclient object and connect to the specified host
\end{cxxdoc}
\end{cxxfunction}
\begin{cxxfunction}
{virtual\ \ }
        {\cxxtilde CBclient}
        {()}
        { Disconnects from the CBserver and deallocates the memory}
        {2.3}
\begin{cxxdoc}

Disconnects from the CBserver and deallocates the memory
\end{cxxdoc}
\end{cxxfunction}
\begin{cxxfunction}
{CBanswer*}
        {tell}
        {(char*\ )}
        { Tells frames to the server}
        {2.4}
\cxxParameter{
\begin{tabular}[t]{lp{0.5\textwidth}}
{\tt\strut char} & *frames the frames
\end{tabular}}
\cxxReturn{
\begin{tabular}[t]{lp{0.5\textwidth}}
{\tt\strut a} & CBanswer object containing the result and the completion
\end{tabular}}
\begin{cxxdoc}

Tells frames to the server


\end{cxxdoc}
\end{cxxfunction}
\begin{cxxfunction}
{CBanswer*}
        {untell}
        {(char*\ )}
        { Untells frames to the server}
        {2.5}
\cxxParameter{
\begin{tabular}[t]{lp{0.5\textwidth}}
{\tt\strut char} & *frames the frames
\end{tabular}}
\cxxReturn{
\begin{tabular}[t]{lp{0.5\textwidth}}
{\tt\strut a} & CBanswer object containing the result and the completion
\end{tabular}}
\begin{cxxdoc}

Untells frames to the server


\end{cxxdoc}
\end{cxxfunction}
\begin{cxxfunction}
{CBanswer*}
        {tellModel}
        {(char**)}
        { Tells files containing frames to the server}
        {2.6}
\cxxParameter{
\begin{tabular}[t]{lp{0.5\textwidth}}
{\tt\strut char**} & files an array of filenames
\end{tabular}}
\cxxReturn{
\begin{tabular}[t]{lp{0.5\textwidth}}
{\tt\strut a} & CBanswer object containing the result and the completion
\end{tabular}}
\begin{cxxdoc}

Tells files containing frames to the server


\end{cxxdoc}
\end{cxxfunction}
\begin{cxxfunction}
{CBanswer*}
        {ask}
        {(char*\ query,\ char*\ format="OBJNAMES",\ char*\ answerrep="FRAME",\ char*\ rollbacktime="Now")}
        { Sends a query to the ConceptBase server}
        {2.7}
\cxxParameter{
\begin{tabular}[t]{lp{0.5\textwidth}}
{\tt\strut char} & *query the query
\\
{\tt\strut char*} & format the format of the query (FRAMES or OBJNAMES)
\\
{\tt\strut char*} & answerrep the format of the answer (FRAME)
\\
{\tt\strut char*} & rollbacktime Rollback Time (e.g.\ "Now")
\end{tabular}}
\cxxReturn{
\begin{tabular}[t]{lp{0.5\textwidth}}
{\tt\strut a} & CBanswer object containing the result and the completion
\end{tabular}}
\begin{cxxdoc}

Sends a query to the ConceptBase server


\end{cxxdoc}
\end{cxxfunction}
\begin{cxxfunction}
{CBanswer*}
        {hypoAsk}
        {(char*\ frames,\ char*\ query,\ char*\ format="OBJNAMES",\ char*\ answerrep="FRAME",\ char*\ rollbacktime="Now")}
        { Sends frames and a query to the ConceptBase server.}
        {2.8}
\cxxParameter{
\begin{tabular}[t]{lp{0.5\textwidth}}
{\tt\strut char} & *frames frames to be told
\\
{\tt\strut char} & *query the query
\\
{\tt\strut char*} & format the format of the query (FRAMES or OBJNAMES)
\\
{\tt\strut char*} & answerrep the format of the answer (FRAME)
\\
{\tt\strut char*} & rollbacktime Rollback Time (e.g.\ "Now")
\end{tabular}}
\cxxReturn{
\begin{tabular}[t]{lp{0.5\textwidth}}
{\tt\strut a} & CBanswer object containing the result and the completion
\end{tabular}}
\begin{cxxdoc}

Sends frames and a query to the ConceptBase server. The frames are told temporarely,
the query is evaluated, and the temporarely objects are removed.


\end{cxxdoc}
\end{cxxfunction}
\begin{cxxfunction}
{CBanswer*}
        {askObjNames}
        {(char*\ query,\ char*\ answerrep="FRAME",\ char*\ rollbacktime="Now")}
        { Sends a query to the ConceptBase server.}
        {2.9}
\cxxParameter{
\begin{tabular}[t]{lp{0.5\textwidth}}
{\tt\strut char} & *query the query
\\
{\tt\strut char*} & answerrep the format of the answer (FRAME)
\\
{\tt\strut char*} & rollbacktime Rollback Time (e.g.\ "Now")
\end{tabular}}
\cxxReturn{
\begin{tabular}[t]{lp{0.5\textwidth}}
{\tt\strut a} & CBanswer object containing the result and the completion
\end{tabular}}
\begin{cxxdoc}

Sends a query to the ConceptBase server. Same as ask but with fixed query format (OBJNAMES).


\end{cxxdoc}
\end{cxxfunction}
\begin{cxxfunction}
{CBanswer*}
        {askFrames}
        {(char*\ query,\ char*\ answerrep="FRAME",\ char*\ rollbacktime="Now")}
        { Sends a query to the ConceptBase server.}
        {2.10}
\cxxParameter{
\begin{tabular}[t]{lp{0.5\textwidth}}
{\tt\strut char} & *query the query
\\
{\tt\strut char*} & answerrep the format of the answer (FRAME)
\\
{\tt\strut char*} & rollbacktime Rollback Time (e.g.\ "Now")
\end{tabular}}
\cxxReturn{
\begin{tabular}[t]{lp{0.5\textwidth}}
{\tt\strut a} & CBanswer object containing the result and the completion
\end{tabular}}
\begin{cxxdoc}

Sends a query to the ConceptBase server. Same as ask but with fixed query format (FRAMES).


\end{cxxdoc}
\end{cxxfunction}
\begin{cxxfunction}
{int}
        {enrollMe}
        {(char*\ host,\ int\ port,\ char*\ user=NULL,\ char*\ tool=NULL)}
        { Connects to a ConceptBase Server Return the return value of connect\_CB\_server (see CBinterfaceh): -1: if socket to specified can not be openend 0: ok other: a completion value (see CBinterfaceh)}
        {2.11}
\cxxParameter{
\begin{tabular}[t]{lp{0.5\textwidth}}
{\tt\strut host} & hostname of the machine where the server runs
\\
{\tt\strut port} & port number of server
\\
{\tt\strut *user} & the name of the tool
\\
{\tt\strut *tool} & the name of the user
\end{tabular}}
\begin{cxxdoc}

Connects to a ConceptBase Server
Return the return value of connect\_CB\_server (see CBinterfaceh):
-1: if socket to specified can not be openend
0: ok
other: a completion value (see CBinterfaceh)


\end{cxxdoc}
\end{cxxfunction}
\begin{cxxfunction}
{int}
        {cancelMe}
        {()}
        { Disconnects from a ConceptBase Server}
        {2.12}
\begin{cxxdoc}

Disconnects from a ConceptBase Server

Return the return value of disconnect\_CB\_server (see CBinterface.h):
-1: error, not connected
0: ok
other: a completion value (see CBinterface.h)
\end{cxxdoc}
\end{cxxfunction}
\begin{cxxfunction}
{CBanswer*}
        {stopServer}
        {(char*\ password=NULL)}
        { Stops the ConceptBase server.}
        {2.13}
\cxxReturn{
\begin{tabular}[t]{lp{0.5\textwidth}}
{\tt\strut a} & CBanswer object containing the result and the completion
\end{tabular}}
\begin{cxxdoc}

Stops the ConceptBase server. Note that a server may be stopped only
by the user who has started it.


\end{cxxdoc}
\end{cxxfunction}
\begin{cxxfunction}
{Clients*}
        {reportClients}
        {()}
        { Return a list of clients connected to the CB server.}
        {2.14}
\begin{cxxdoc}

Return a list of clients connected to the CB server.
The result will be a list of Client objects as defined in libCB.
\end{cxxdoc}
\end{cxxfunction}
\begin{cxxfunction}
{CBanswer*}
        {nextMessage}
        {(char*\ method="")}
        { Gets a message from the server}
        {2.15}
\cxxParameter{
\begin{tabular}[t]{lp{0.5\textwidth}}
{\tt\strut char*} & method the type of the message to be retrieved
\end{tabular}}
\cxxReturn{
\begin{tabular}[t]{lp{0.5\textwidth}}
{\tt\strut a} & CBanswer object containing the result and the completion
\end{tabular}}
\begin{cxxdoc}

Gets a message from the server


\end{cxxdoc}
\end{cxxfunction}
\begin{cxxfunction}
{CBerror*}
        {getErrorMessages}
        {()}
        { Gets the error messages from the server}
        {2.16}
\cxxReturn{
\begin{tabular}[t]{lp{0.5\textwidth}}
{\tt\strut a} & string containing all error messages
\end{tabular}}
\begin{cxxdoc}

Gets the error messages from the server


\end{cxxdoc}
\end{cxxfunction}
\begin{cxxfunction}
{CBanswer*}
        {LPICall}
        {(char*\ )}
        { Perform a LPI call on the server.}
        {2.17}
\begin{cxxdoc}

Perform a LPI call on the server. A LPI call is a call of Prolog-predicate
of the CBserver. This is mostly used for debugging.
\end{cxxdoc}
\end{cxxfunction}
\begin{cxxfunction}
{inline\ \ \ int}
        {connected}
        {()}
        { Check whether this client is connected}
        {2.18}
\begin{cxxdoc}

Check whether this client is connected
\end{cxxdoc}
\end{cxxfunction}
\begin{cxxfunction}
{inline\ \ }
        {operator\ int}
        {()}
        { The operator int checks also if the client is connected.}
        {2.19}
\begin{cxxdoc}

The operator int checks also if the client is connected.

\end{cxxdoc}
\end{cxxfunction}
\begin{cxxfunction}
{char*}
        {getServerName}
        {()}
        { Return the name of the server}
        {2.20}
\begin{cxxdoc}

Return the name of the server
\end{cxxdoc}
\end{cxxfunction}
\begin{cxxfunction}
{char*}
        {getClientName}
        {()}
        { Return the name of the client}
        {2.21}
\begin{cxxdoc}

Return the name of the client
\end{cxxdoc}
\end{cxxfunction}
\end{cxxclass}


\begin{cxxclass}
{class}
        {CBanswer}
        {}
        { C++ Wrapper for Answer struct of libCB}
        {1}
\begin{cxxfunction}
{}
        {CBanswer}
        {(Answer*\ ans)}
        {Constructs a CBanswer object from a Answer struct }
        {1.1}
\cxxParameter{
\begin{tabular}[t]{lp{0.5\textwidth}}
{\tt\strut ans} & pointer to the Answer struct
\end{tabular}}
\begin{cxxdoc}
Constructs a CBanswer object from a Answer struct

\end{cxxdoc}
\end{cxxfunction}
\begin{cxxfunction}
{}
        {\cxxtilde CBanswer}
        {()}
        {Deallocate the memory of the object }
        {1.2}
\begin{cxxdoc}
Deallocate the memory of the object
\end{cxxdoc}
\end{cxxfunction}
\begin{cxxfunction}
{Completion}
        {getCompletion}
        {()}
        {Get the completion value of the answer }
        {1.3}
\begin{cxxdoc}
Get the completion value of the answer
\end{cxxdoc}
\end{cxxfunction}
\begin{cxxfunction}
{char*}
        {getResult}
        {()}
        {Get the result string of the answer }
        {1.4}
\begin{cxxdoc}
Get the result string of the answer
\end{cxxdoc}
\end{cxxfunction}
\begin{cxxfunction}
{char*}
        {getRespondingTool}
        {()}
        {Get the ID of the responding tool of the answer.}
        {1.5}
\begin{cxxdoc}
Get the ID of the responding tool of the answer. This
usually the CBserver
\end{cxxdoc}
\end{cxxfunction}
\end{cxxclass}




\begin{cxxclass}
{class}
        {CBerror}
        {}
        { C++ Wrapper of Error\_Messages struct in libCB.}
        {3}
\begin{cxxfunction}
{}
        {CBerror}
        {(Error\_Messages*\ e)}
        {Construct a CBerror object from a list of Error\_Messages }
        {3.1}
\cxxParameter{
\begin{tabular}[t]{lp{0.5\textwidth}}
{\tt\strut e} & pointer to the Error\_Messages
\end{tabular}}
\begin{cxxdoc}
Construct a CBerror object from a list of Error\_Messages

\end{cxxdoc}
\end{cxxfunction}
\begin{cxxfunction}
{}
        {\cxxtilde CBerror}
        {()}
        { Deallocate the memory of a CBerror object}
        {3.2}
\begin{cxxdoc}

Deallocate the memory of a CBerror object
\end{cxxdoc}
\end{cxxfunction}
\begin{cxxfunction}
{char*}
        {getErrorMessage}
        {()}
        { Get the error message of this object.}
        {3.3}
\begin{cxxdoc}

Get the error message of this object. This will return only
the first error message of the list.
\end{cxxdoc}
\end{cxxfunction}
\begin{cxxfunction}
{char*}
        {getAllErrorMessages}
        {()}
        { This method will return all error messages of the list.}
        {3.4}
\begin{cxxdoc}

This method will return all error messages of the list. The method
will allocate a new string, thus the resulting string has to be freed
by the calller.
\end{cxxdoc}
\end{cxxfunction}
\begin{cxxfunction}
{CBerror*}
        {getNextError}
        {()}
        { Get the next error message in the list}
        {3.5}
\begin{cxxdoc}

Get the next error message in the list
\end{cxxdoc}
\end{cxxfunction}
\end{cxxclass}
